\subsection{A ceiling for the homotopy}\label{e:ceiling-section}

We keep the notation and the four stages of Section~\ref{d:homotopy}.
All integrals without a stated measure use the metric $g$; $\nu$ at an
inner face points into the retained exterior.  A constant $\Pi_n$ is
bounded by a polynomial in $n$, with the prepared data, collars and
cutoff shapes fixed.  In contrast, $C(n)$ may be arbitrary at fixed
$(n,\eps)$; it is always independent of the sufficiently large
truncation radius and of the homotopy parameter.

\begin{lemma}[Weighted trace estimates]\label{e:weighted-trace}
In the first two homotopy stages, for every nonnegative smooth
$\varphi$ one has
\begin{align}
 \int_B L\varphi
 &\le \Pi_n\int_{\mathrm{coll}}u
 \bigl[(1+\sqrt{\mathcal T})\varphi+
                         |\dd\varphi|_{A_\chi}\bigr],
 \label{e:weighted-trace-eq}\\
 \int_B \varphi
 &\le \Pi_n\int_{\mathrm{coll}}\sqrt D
 \bigl[(1+\sqrt{\mathcal T})\varphi+
                         |\dd\varphi|_{A_\chi}\bigr].
 \label{e:unweighted-trace}
\end{align}
The integrands on the right include fixed collar cutoffs.  In the first
stage, moreover,
\begin{equation}\label{e:small-boundary-flux}
 |(V_\lambda)_\nu|\le Cud
                 \le C\frac{\tau_0}{\ell}L.
\end{equation}
None of these estimates requires an upper bound for $Z$.
\end{lemma}

\begin{proof}
Set $W=w/\sqrt D$.  By the two-sided normal-slope bounds
\eqref{d:collar-slopes}, $w\cdot\nu$ has a definite sign on each face and
$a\ge1/2$ there.  The corresponding signed fluxes of $uW=Lw$ and
$\sqrt D W=w$ are $La$ and $a$.  Apply the divergence theorem to these
fields, multiplied by the collar cutoff and $\varphi$, reversing the
sign on the white collars.  The estimates needed for the resulting
volume terms are
\begin{gather}
 |W(\varphi)|\le |\dd\varphi|_{A_d}
                    \le \Pi_n|\dd\varphi|_{A_\chi},
 \label{e:trace-direction}\\
 \frac{|\divg(uW)|}{u}
       +\frac{|\divg(\sqrt D W)|}{\sqrt D}
                   \le\Pi_n(1+\sqrt{\mathcal T}).
 \label{e:trace-divergence}
\end{gather}
For the second inequality, direct differentiation gives
\[
 \divg w=\tr_{A_d}H^f+d\,p\,w(t).
\]
In $\divg(Lw)/u$ every derivative is multiplied by $\sqrt d$;
in particular the longitudinal Hessian contribution is controlled by
$\sqrt d\,|q|$.  Derivatives of $L$ contribute
$(p+2)\sqrt d\,w(t)$.  These and the transverse terms are bounded by
\eqref{d:coercivity}; replacing $H^f$ by $S_1-K$ costs a bounded
compactly supported term.  The same calculation applies to the second
quotient in \eqref{e:trace-divergence}.  Derivatives of the fixed collar
cutoffs have bounded cost.  This proves both trace estimates.

At assigned first-stage heights, the two colors both satisfy
$w_\nu(F-P_B)=aH$.  Consequently the physical boundary expression is
\[
 \mathcal B=(a-1)H-N_Bd
      =-d\left(N_B+\frac{H}{1+a}\right).
\]
The omitted normal drift in the first homotopy stage has size
$O(\chi)\le O(d)$, because $\dd f$ is normal on a face and $A_\chi$
has normal eigenvalue $\chi$.  The prefactor in the homotopy boundary
condition lies in $[0,1]$.  Thus its flux has size $O(ud)$.
Finally, \eqref{d:collar-slopes} and $l\ge\ell$ give
$D^{-1/2}\le C\tau_0/\ell$, proving
\eqref{e:small-boundary-flux}.
\end{proof}

\begin{proposition}[Ceiling and bounded variable range]\label{e:ceiling}
For all sufficiently large $n$, and then all sufficiently large
truncations $\Omega_R$, every smooth homotopy solution with
$t\ge-\eps$ satisfies
\[
 |Z|+|t|+|f|+|\nabla f|\le C(n).
\]
The constants are uniform over the four stages and $R\ge R_{\min}(n)$.
\end{proposition}

\begin{proof}
Only an upper bound for $Z$ remains to be proved.  Indeed
$Z=t+\tfrac18\log D\ge-\eps$, and an upper bound for $Z$, together
with $t\ge-\eps$ and $l\ge\ell$, bounds $D$, $t$ and $\sigma$.
The height estimate already bounds $f=h/\tau_0$ at fixed $n$.

\paragraph{An integral estimate above a fixed height.}
There is $M_0=M_0(n)$ such that, in the first three stages,
\begin{equation}\label{e:high-source}
 Z>M_0\quad\Longrightarrow\quad
 \Xi\ge\delta_0\mathcal T+\rho+
                         \frac{\delta_0\tau_0}{2}a\sigma.
\end{equation}
In fact the penalty is supported, on the allowed side of the floor,
where $-\eps\le t\le-\eps+1/n$.  There $l$ and $L$ are
comparable to $\ell$, so large $Z$ forces arbitrarily large $\sigma$
and hence arbitrarily large $a\sigma$.  The bounded coefficient
$\mathcal P=C_n(\rho_0+v\rho_1)$ is therefore absorbed into half of
the indicated height term.

In the first stage choose a smooth increasing $k_0$ which vanishes
below $M_0$ and equals one above $M_0+1$, and test the divergence
equation with $k_0(Z)$.  Integration by parts gives the nonnegative
principal contribution
$4\int u k_0'|\dd Z|_{A_\chi}^2$ on the source side.  The drift cross
term is bounded by half of this contribution plus
$C\int u k_0'|K|^2$.  The latter is bounded: it has compact support,
and in the transition strip $Z\in[M_0,M_0+1]$ the floor bounds $D$, $u$
and all zeroth-order coefficients.
The inner flux costs at most
\[
 C\frac{\tau_0}{\ell}\int_B Lk_0.
\]
Use \eqref{e:weighted-trace-eq}.  On the fixed collars, positivity of
$\rho$ implies
$1+\sqrt{\mathcal T}\le C(\rho+\delta_0\mathcal T)$, and
$\Pi_n\tau_0/\ell\to0$.  Thus these terms absorb into the positive
source in \eqref{e:high-source}.  The remaining trace term contains
$k_0'|\dd Z|_{A_\chi}$; Young's inequality absorbs its quadratic part
into the principal contribution, with another bounded
transition-strip remainder.  We obtain
\begin{equation}\label{e:high-integral}
 \int_{\Omega_R}u k_0(Z)
          (\mathcal T+\rho+\tau_0a\sigma)\le C(n).
\end{equation}

Let $\Gamma$ denote the ordinary graph of $f$ in $g+\dd b^2$, with
measure $\dd\Gamma=\sqrt{1+\sigma^2}\,\dd V_g$.  On every fixed
compact set, \eqref{e:high-integral} bounds $e^Z$ in
$L^2(\dd\Gamma)$.  To see the density comparison, use
\[
 e^{2Z}=e^{2t}D^{1/4},\qquad
 D^{1/4}\le C(n)(1+\tau_0a\sigma),\qquad
 \sqrt{1+\sigma^2}\asymp_n\sqrt D.
\]
The part where $k_0=1$ is controlled by
\eqref{e:high-integral}, since $\rho$ has a positive lower bound on
that compact set.  On the complementary part $Z\le M_0+1$ the floor
bounds the entire density.  Moreover $\ell\le l\le\exp(1)$ and
$d/\chi\le1+n/4$, so the graph gradient norm is comparable, at fixed
$n$, to $|\dd\cdot|_{A_\chi}$.

\paragraph{Sobolev inequality on the graph, including its boundary.}
On compact base patches, including patches meeting an inner face,
\begin{equation}\label{e:graph-sobolev}
 \|\omega\|_{L^6(\dd\Gamma)}^2
 \le C(n)\int_\Gamma
       \bigl(|\nabla_\Gamma\omega|^2+
                         (1+\mathcal T)\omega^2\bigr)
\end{equation}
for smooth $\omega$ supported away from the other patch edges.
Here the constant is independent of the particular graph.
Embed a fixed smooth extension of the base patch isometrically in a
Euclidean space, using Nash's theorem \cite{Nash1956}, and use the
product embedding for the graph.  The Euclidean mean-curvature vector
of this graph obeys
\[
 |\mathbf H_\Gamma|\le C(n)(1+\sqrt{\mathcal T}).
\]
Indeed its graph-normal component is the trace of
$\nabla^2f/\sqrt{1+\sigma^2}$ with longitudinal coefficient
$(1+\sigma^2)^{-1}$; all its transverse entries and this longitudinal
trace are controlled by \eqref{d:coercivity}.  The second fundamental
form of the fixed base embedding adds a bounded term, since it is
contracted against the inverse graph metric.

The Euclidean submanifold Sobolev inequality of Michael--Simon
\cite{MichaelSimon1973}, in the boundary form of
\cite[Theorem~1]{Brendle2021}, therefore gives
\[
 \left(\int_\Gamma\varphi^{3/2}\right)^{2/3}
 \le C(n)\left[
 \int_\Gamma\bigl(|\nabla_\Gamma\varphi|+
                   (1+\sqrt{\mathcal T})\varphi\bigr)
                  +\int_{\partial\Gamma}\varphi\right]
\]
for nonnegative $\varphi$ of the stated support.  The graph has
constant height on each inner face, so its boundary measure is exactly
$\dd A_g$.  The last term is controlled by
\eqref{e:unweighted-trace} and the graph comparisons above.  Taking
$\varphi=|\omega|^4$, then applying Cauchy--Schwarz to
$|\omega|^3|\nabla_\Gamma\omega|$ and
$(1+\sqrt{\mathcal T})|\omega|^4$, proves
\eqref{e:graph-sobolev}.

\paragraph{Iteration to a supremum bound.}
For a fixed base cutoff $\eta$ test the first-stage equation with
$\eta^2e^{2kZ}/L$, where $k\ge k_*(n)$ will be large.  After integration
by parts the positive differentiated-test contribution is
\[
 8k\int \frac{u}{L}\eta^2e^{2kZ}|\dd Z|_{A_\chi}^2.
\]
The source has lower bound $\delta_0\mathcal T-\Pi_n$.
Since $\dd\log L=(p+2)\dd t$, coercivity and Young's inequality bound
its cross terms with $4\dd Z+\lambda K(w,\cdot)$ by
\[
 \frac{\delta_0}{2}\mathcal T+
                  C(n)(1+|\dd Z|_{A_\chi}^2).
\]
Take $k_*(n)$ large enough to absorb the gradient-square part;
control drift and cutoff terms by the same inequality.  After
cancelling $L$, the boundary term is bounded by
$C\int_B\eta^2e^{2kZ}$.  Apply \eqref{e:unweighted-trace} to
$\varphi=\eta^2e^{2kZ}$.  Its $\sqrt{\mathcal T}$ contribution is
absorbed into a small part of the source.  Its differentiated
exponential contributes $C(n)k|\nabla_\Gamma Z|$, which costs a
small part of $k|\nabla_\Gamma Z|^2$ and $C(n)k$ by Young's
inequality.  Thus
\begin{equation}\label{e:moser-energy}
 \int_\Gamma e^{2kZ}\eta^2
          (\mathcal T+k|\nabla_\Gamma Z|^2)
 \le C(n)k\int_\Gamma e^{2kZ}
                              (\eta^2+|\dd\eta|_g^2).
\end{equation}
Applying \eqref{e:graph-sobolev} to $\omega=\eta e^{kZ}$ and using
\eqref{e:moser-energy} gives the iteration from exponent $2k$ to
$6k$.  For nested patches with separation $s$, the multiplier in
this step is at most
\[
 [C(n)k^2(1+s^{-2})]^{1/(2k)}.
\]
The product of these multipliers converges when $k$ is multiplied by
three and the patch gaps decrease geometrically.  Hence on two nested
patches $U'\Subset U$ (relative to the inner boundary),
\[
 \sup_{U'}e^Z\le C(n,s)
                  \|e^Z\|_{L^{2k_*}(U,\dd\Gamma)},
\]
with $C(n,s)$ bounded by a fixed inverse power of $s$.
Interpolation with the already bounded $L^2$ norm gives
\[
 \sup_{U'}e^Z\le C(n,s)
             (\sup_Ue^Z)^{1-1/k_*}.
\]
Young's inequality, with any desired small coefficient for
$\sup_Ue^Z$, and a second iteration on geometrically nested patches
absorb the larger supremum.  The terminal term vanishes because each
individual smooth solution has a finite supremum; the resulting bound
is uniform over solutions.  This proves $Z\le C(n)$ on the compact
sets containing the drift support and all inner faces.

Outside the drift support, a high interior maximum contradicts
\eqref{e:high-source} by the divergence equation.  Comparison, with
zero outer data and the compact bounds just proved, completes the
first-stage ceiling.  In the second and third stages the drift is zero
everywhere, so this maximum argument applies directly.  At a boundary
maximum with $Z$ sufficiently large, the upper normal-slope bound of
Proposition~\ref{d:collars} bounds $D$ on the face, hence forces
$t>1$.  Then $j(t)=1$, so the prescribed conormal flux is zero.
The Hopf boundary lemma excludes such a high maximum.  The same
argument applies in the fourth stage, where $f=0$ and the source and
flux are scaled together.  At scale zero the mixed homogeneous
problem has only the zero solution.

The required choices of $n$ and $R_{\min}(n)$ are compatible with
Proposition~\ref{d:collars}.  More explicitly, failure of a uniform
positive height gap for all $n\ge N$ and $R\ge R_{\min}(n)$,
for every choice of $N$ and finite $R_{\min}$, would give, for each
$j$, an $n_j\ge j$ admitting arbitrarily large failing radii at gap
$1/j$.  Choose such an $R_j\ge j$.  This is precisely the contrary
sequence excluded in Lemma~\ref{d:height-separation}.  A finite
number of exceptional fixed $n$ can be discarded by increasing $N$;
infinitely many with unbounded failure tails would give the same
contrary sequence.  After fixing the uniform gap, eventual validity
at each remaining $n$ defines a finite $R_{\min}(n)$, which we may
enlarge to be at least $n$.  No uniform estimate in $n$ is asserted
for the subsequent ceiling constant.
\end{proof}
